\chapter{Present Obligations and Distant Futures}
\label{ch:31}

This chapter addresses the feature of the program that most closely resembles the ideologies the interview criticizes and that therefore requires the most careful treatment. Longtermism, as the interview presents it, is the view that the welfare of future persons deserves equal weight with the welfare of the living, and that if harms to present people are required in order to secure the welfare of a vastly larger number of future beings, those harms are acceptable. The interview presents this as a rationale for sacrifice in the present on behalf of speculative future populations and as a component of a larger set of commitments in which the construction of a successor intelligence is justified by the scale of its imagined benefits. This book need not endorse the interview's account of the history or the personnel of these movements to recognize the logical structure at issue. The target of the argument is a particular decision rule, namely unconstrained maximization of expected total welfare over all future persons without discount, and not longtermism as a family of positions, many of which do not endorse that rule. When the number of future beneficiaries is allowed to be astronomically large and their welfare is weighted without discount, the arithmetic of unconstrained maximization is dominated by the future, and any present cost, however severe, is justified if it increases the probability of the future benefit by even a tiny amount. The reasoning is vulnerable on several grounds: its conclusions depend on probability estimates that cannot be verified, it is sensitive to modeling assumptions that are not independently checkable, and it invites exactly the kind of justification that is difficult for those who bear the cost to contest.

The program responds not by assigning the future zero weight, which would be arbitrary, but by changing the structure of the decision. Obligations to existing persons are treated as constraints that define the set within which optimization over the future takes place, instead of as terms in a sum that can be outweighed. A person alive now is owed minimum provision of the conditions of a decent life, together with the procedural rights through which they can contest decisions affecting them, and these entitlements are not tradeable against benefits to others, whether the others are alive or not yet born. Within the set of options that respect the floor, the long-term consequences of choices are given substantial weight, and the program regards the preservation of the conditions under which future persons can live well as a genuine obligation. The constraint is also a response to uncertainty: because the estimates on which speculative benefits rest are unverifiable, decisions that affect the present should be robust across the range of plausible models rather than dependent on any one. This is the same demand for independent verification applied in Chapter 9, now applied to the justification of sacrifice.

The distinction thereby drawn between this program and the ambitions the interview describes is institutional and does not rest on a claim about intentions. The program holds that capability be directed by procedures answerable to living persons, that the costs and benefits of projects be allocated by processes those persons can contest, that claims of benefit be verified by parties independent of those who make them, that the delegation of authority to machines be limited by the principle stated at the outset, and that present obligations constitute a floor and not a price. A project that meets these conditions can include astronomical engineering and extreme longevity without being a project whose legitimacy is established by the grandeur of its goal. A project that fails them does not become acceptable by being benevolently intended.

The argument from the weight of the future is usually made in the language of discounting, and the language is worth examining, because the disagreement it conceals is instructive. In the standard framework, associated with Ramsey's analysis of saving in 1928~\cite{ramsey1928}, the value of a stream of future welfare is computed by weighting each period by a discount factor, composed of a pure rate of time preference, which gives less weight to the future simply because it is future, and a term reflecting the growth of consumption and the diminishing marginal value of additional consumption. The Stern Review of 2006 set the pure rate near one-tenth of one per cent a year, on the ground that to prefer the present to the future merely because of the date is without ethical justification, and obtained a high value for the prevention of climate harm; critics, among them Nordhaus, argued for rates an order of magnitude higher, partly on the ground that they reflect the observed behavior of savers and investors. The dispute does not turn on any fact that further research might discover, since it is a dispute about how to weigh persons by when they live, and the program's response is to avoid making the conclusion depend on it. A floor of provision for the living that is not subject to discounting, together with optimization over the remaining set, makes the choice of rate relevant only within the admissible set, where its influence is limited, and it is the use of the rate to override the floor that the program refuses.

A second difficulty is the uncertainty of the long future. The expected value of an action that affects the very distant future is the sum of a large number of terms, each of which is the product of a very small probability and a very large payoff, and the terms depend on estimates which cannot be checked. The result is a form of argument in which a speculative probability of benefit, multiplied by an astronomical number of beneficiaries, outweighs any certain cost to the living, and in which the conclusion can be driven in either direction by adjustments of the probability that no available evidence constrains. This is the structure that has been called Pascal's mugging, after an example in which a stranger offers an enormous payoff for a small payment on a claim that cannot be verified, and it is the form in which the argument from the future becomes unfalsifiable. The remedy proposed in this book is the one that has been used throughout, the demand for independent verification, applied here to the claims on which a sacrifice would rest: a proposal to impose a cost on present persons for the sake of future ones must supply an estimate of the probabilities involved, the evidence for them, the sensitivity of the conclusion to the estimate, and a statement of what would count as evidence against it. A proposal that cannot do this has not stated a claim that can bind anyone.

A third consideration is that obligations to future persons are not obligations to a determinate set of individuals. Parfit's non-identity problem observes that the policies we adopt determine which persons will exist, so that a policy cannot be said to harm a particular future person who would not have existed had a different policy been adopted. The observation does not show that there are no duties to the future, since the existence of persons whose lives are bad is itself something that can be wrong, but it complicates the application of duties of harm to future persons and favors formulations in terms of the conditions that any future person should be able to expect. A principle of this kind, which Rawls put forward as a principle of just saving~\cite{rawls1971}, holds that each generation should pass on to the next conditions at least as favorable as those it inherited, in respect of the institutions that sustain justice and the stock of capital, natural and made, on which a decent life depends. The program adopts this as a minimal intergenerational floor, complementary to the floor for the living, and it does not require that the future be treated as a source of claims that outweigh those of the present.

What does the floor for the living contain, and how is it connected to the readiness classification of Chapter 16? It contains clean water and sanitation, the energy needed for a decent life, health care, education, and the political and legal rights through which a person can contest what is done to them. These are all matters on which the technologies of the first readiness class are available or nearly so, and the programs of the megascale parts need not be completed, or even undertaken, in order to meet the floor. The rule that no measure that satisfies a present obligation may depend on a program of the second or third class is the means of keeping the two apart, and the audit of the plan is the means of enforcing it. The audit has a second function, which is to make the cost of the floor visible. The loss of potential value that the floor imposes, defined in the formalization, is a number that can be estimated, however roughly, and its publication allows the persons who bear the cost of the floor, and those who would benefit from its relaxation, to see what is being traded. A society that chooses to maintain the floor has made a choice, and a society that does not know the price of its choice has not made one.

The relation of this chapter to the preceding parts should be stated, because the chapters on space and on longevity are the places where the argument from the future has the greatest force. The reduction of the hazard of extinction by dispersal to other worlds, the preservation of options for the very long term, and the cultivation of knowledge that may be of value to persons in the distant future are defensible aims, and the program does not oppose them. It opposes only their use as grounds for diminishing the provision made for living persons, and as grounds for removing the authority of living persons over the course of events. The distinction is between a program that seeks to leave the future more options and a program that seeks to choose the future's course. The first is compatible with the principle that capability is to be judged by the range of admissible futures it makes available; the second is the displacement of authority that the book was written to resist.

\section*{Formalization}

Let time be indexed by $t\in\{0,1,2,\dots\}$, let $N_t$ be the number of persons alive at $t$, and let $u_{i,t}\in\R$ be the welfare of person $i$ at $t$. An action (policy) $a$ in a set $\mathcal{A}$ induces welfare profiles $\{u_{i,t}(a)\}$. Unconstrained undiscounted utilitarian maximization selects $a\in\arg\max_{a}\ W(a)$ with
\[
W(a)=\sum_{t\ge0}\sum_{i=1}^{N_t}u_{i,t}(a).
\]

\begin{proposition}[Domination by the future]
Let $a'$ be a policy that preserves present welfare and let $a$ be a policy that costs present persons a finite total welfare $\ell>0$ while raising by $\Delta p>0$ the probability of attaining a future of total value $V_f$, other effects being ignored. Then the unconstrained expected-value criterion prefers $a$ to $a'$ if and only if $V_f>\ell/\Delta p$. Consequently, for every finite $\ell>0$ and every $\Delta p>0$ there exists a threshold $V_0=\ell/\Delta p$ such that the criterion favors $a$ whenever the claimed future value exceeds $V_0$. Moreover, if $N_t\to\infty$ with mean welfare bounded below by some $\epsilon>0$, then $W(a)$ diverges and the comparison of policies is undefined without a convention such as discounting or truncation.
\end{proposition}

\begin{proof}
The expected change in $W$ from adopting $a$ is $\Delta p\,V_f-\ell$, which is positive exactly when $V_f>\ell/\Delta p$. For the final claim, $\sum_{t}N_t\bar u_t\ge\epsilon\sum_tN_t=\infty$ when $N_t\to\infty$, where $\bar u_t$ is the mean welfare at time $t$.
\end{proof}

The proposition concerns a decision rule, namely unconstrained maximization of expected total welfare over all future persons, and not longtermism as a family of positions. Many positions described as longtermist do not endorse this rule, and the argument does not apply to them.

The constrained program replaces this with a feasible set. Let $w_{\min}\in\R$ denote the floor of welfare owed to each living person and let $\mathcal{R}$ be the set of procedural rights. Define
\[
\mathcal{A}_{\mathrm{adm}}=\Big\{a\in\mathcal{A}:\ u_{i,0}(a)\ge w_{\min}\ \text{for all } i\le N_0,\ \text{and } a\ \text{respects}\ \mathcal{R}\Big\}.
\]
Let $\mathcal{M}$ be a set of plausible models of the consequences of policies, each assigning to $a$ a value $W_m(a)$. The robust criterion selects
\[
a^\ast\in\arg\max_{a\in\mathcal{A}_{\mathrm{adm}}}\ \min_{m\in\mathcal{M}}\ W_m(a),
\]
which is well-defined when $W_m$ is finite for each $m$, for instance under a cutoff horizon $T_{\max}$ chosen by the legitimate procedure $\Sigma$.

\begin{proposition}
For every $a\in\mathcal{A}\setminus\mathcal{A}_{\mathrm{adm}}$ and every $a'\in\mathcal{A}_{\mathrm{adm}}$, the constrained criterion prefers $a'$ regardless of the magnitude of $W_m(a)-W_m(a')$. In particular no future benefit, however large, justifies a policy that violates the floor or the rights of persons alive at $t=0$.
\end{proposition}

\begin{proof}
The criterion maximizes over $\mathcal{A}_{\mathrm{adm}}$ only, so policies outside the set are never selected, and the selection among members of the set does not depend on the values of policies outside it.
\end{proof}

The proposition formalizes the statement that present obligations are a floor and not a price. Its price is a loss of potential value, namely $\sup_{a\in\mathcal{A}}\min_mW_m(a)-\sup_{a\in\mathcal{A}_{\mathrm{adm}}}\min_mW_m(a)\ge0$, which the program accepts and which should be reported to those who decide.
